\section{Two caches, and what they actually save}
\label{sec:ch03-caches}
\index{caching|(}

Positions four and eight are caches, and they cache different things for
different reasons. The document cache remembers documents the server has already
accepted. The operation cache remembers what compiling one produced. Send the
same query twice and both report a hit:

\begin{minted}[fontsize=\footnotesize]{text}
parse - validate 0.204ms compile 0.097ms coerce - execute 0.492ms total 0.931ms
    (document cache miss, operation cache miss, 146 resolvers)
parse - validate - compile - coerce - execute 1.049ms total 1.065ms
    (document cache hit, operation cache hit, 146 resolvers)
\end{minted}

The dashes on the second line are the point. Validation and compilation did not
run slowly; they did not run.

\subsection*{The document cache is not a parse cache}

\index{caching!validation skipped}%
Calling position four a parse cache is the natural reading of the name and it
undersells what it does. On a hit, \code{DocumentCacheMiddleware} copies six
things onto the request, and one of them is the interesting one:

\begin{minted}{csharp}
documentInfo.Id = request.DocumentId;
documentInfo.Hash = document.Hash;
documentInfo.Document = document.Body;
documentInfo.IsValidated = true;
documentInfo.IsCached = true;
documentInfo.IsPersisted = document.IsPersisted;
\end{minted}

\code{IsValidated}. Two middleware later, that flag is the whole of the
validation decision:

\begin{minted}{csharp}
if (!documentInfo.IsValidated || _documentValidator.HasNonCacheableRules)
{
    using (_diagnosticEvents.ValidateDocument(context))
    {
        ...
    }
}
\end{minted}

So the document cache is better described as a record of documents this schema
has already agreed to run. That is a much stronger claim than remembering a
syntax tree, and it is why the entry is written on the way out rather than on
the way in. \code{DocumentCacheMiddleware} calls \code{TryAddDocument} after
\code{await \_next(context)} returns, and only when \code{IsValidated} is still
true. A document that fails validation is never cached, which is the behaviour
you want: caching it would mean caching permission to run something the schema
rejects.

\index{validation!non-cacheable rules}%
The second half of that condition is the escape hatch, and I want to be careful
about what I know. If the validator holds rules that cannot be cached, every
request is validated again whatever the cache says. Mosaic's measurements show
validation being skipped, so its configuration has no such rules; which of
HotChocolate's shipped rules are non-cacheable I did not establish, and I am not
going to guess in a chapter whose whole argument is that you can go and look. If
you add custom validation rules, check this branch before assuming the cache
still saves you the work.

\subsection*{The operation cache, and the thundering herd}

\index{caching!operation cache}%
Position eight caches the compiled operation, and its implementation is more
interesting than a dictionary lookup. Consider a cold start behind a load
balancer: fifty requests for the same uncached query arrive in the same
millisecond. A naive cache lets all fifty miss and all fifty compile.

\code{OperationCacheMiddleware} does not. It keeps a second dictionary of
compilations currently in flight, and the first request to claim an entry
becomes what the code calls the single-flight leader:

\begin{minted}{csharp}
if (ReferenceEquals(cachedInFlightOperation, inFlightOperation))
{
    // We won the race! This request is the single-flight leader
    // responsible for compiling and signaling all followers.
    isSingleFlightLeader = true;
    context.Features.Set(inFlightOperation.Value);
}
else
{
    // We lost the race! Another request claimed leadership between
    // TryGetValue and GetOrAdd. So we simply await the leader's result.
    var coalescedOperation = await cachedInFlightOperation.Value.Task
        .WaitAsync(context.RequestAborted)
        .ConfigureAwait(false);
    context.SetOperation(coalescedOperation);
}
\end{minted}

Those comments are HotChocolate's, exclamation marks and all. The leader
compiles, the followers wait on its result, and the leader adds the operation to
the cache before removing the in-flight entry so there is never a moment when it
is in neither. I have not measured this path, because manufacturing the race
reliably is a chapter of its own, but it is worth knowing it exists: it is the
difference between a cold restart costing one compilation and costing as many as
you have concurrent clients.

\subsection*{What the caches are worth, honestly}

\index{caching!measuring}%
Here is where I nearly published a number I would have had to retract.

The first request to a freshly started Mosaic reports this:

\begin{minted}[fontsize=\footnotesize]{text}
parse - validate 16.690ms compile 11.515ms coerce - execute 41.541ms total 88.061ms
\end{minted}

Eighty-eight milliseconds, against one or two for the same query once the caches
are warm. It is a wonderful statistic and it is mostly a lie. Almost all of it
is the runtime compiling itself, because the validator, the operation compiler
and the resolver machinery are all being jitted for the first time. Quote it as
the cost of validation and you have told your reader that validating a
nine-field query takes sixteen milliseconds, which it does not.

The way to get an honest number is to warm the process first and then feed it
work it has never seen. I ran the catalog-with-reviews query twenty times to
settle the runtime down, then sent five documents that differed only by an
alias, so each one was novel to both caches while doing identical work. On this
machine, those five reported:

\begin{itemize}
  \item validation between 0.184 and 0.355~ms,
  \item compilation between 0.076 and 0.110~ms,
  \item and a total inside the pipeline between 0.75 and 1.33~ms.
\end{itemize}

Two things follow. The first is that the caches save something like a third of a
warm request for this query, which is worth having and is not the order of
magnitude the cold number implied. The second is more useful: validation costs
roughly three times what compilation costs. If you had asked me to guess which
of those two was the expensive one, I would have said the one called
\enquote{compile}, and I would have been wrong by a factor of three.

The absolute milliseconds are worth nothing on your machine. The ratio, and the
habit of warming a process before believing a first-request measurement, are
worth carrying. What the caches hand on once they have let a document through is
the subject of the next section.
\index{caching|)}
